\section{System Architecture}
\label{sec:architecture}

QuantumCraft is built on Fabric, the same modding toolchain CircuitSimCraft
uses~\cite{picos2026circuitsimcraft}. The codebase is organized into five packages, shown in
Figure~\ref{fig:architecture}: a simulator package with no dependency on any Minecraft class
at all, a pair of packages describing the in-world blocks and their per-block state, a
networking package that scans a player's placed blocks into a circuit and reports the result
back to the client, and a client-only package responsible for rendering the probe's result
heads-up display.

\begin{figure}[htbp]
\centering
\begin{tikzpicture}[
    box/.style={draw, rounded corners, minimum width=6.4cm, minimum height=1cm,
                align=center, font=\small},
    node distance=0.45cm
]
\node[box, fill=gray!10] (sim) {\texttt{sim} --- pure-Java statevector simulator\\(no Minecraft dependency)};
\node[box, below=of sim, fill=blue!8] (blockentity) {\texttt{blockentity} --- qubit-line, gate, and\\control/target marker state};
\node[box, below=of blockentity, fill=blue!8] (block) {\texttt{block} --- placement and right-click preset cycling};
\node[box, below=of block, fill=orange!12] (network) {\texttt{network} --- world-scan-to-circuit bridge,\\probe result protocol};
\node[box, below=of network, fill=green!10] (client) {\texttt{client} --- result HUD rendering (client-side only)};

\draw[-{Latex[length=2mm]}] (blockentity) -- (sim);
\draw[-{Latex[length=2mm]}] (block) -- (blockentity);
\draw[-{Latex[length=2mm]}] (network) -- (blockentity);
\draw[-{Latex[length=2mm]}] (network) -- (sim);
\draw[-{Latex[length=2mm]}] (client) -- (network);
\end{tikzpicture}
\caption{Package-level architecture. As in CircuitSimCraft, the simulator has zero Minecraft
dependencies and is directly unit-testable with plain \texttt{javac}/\texttt{java}; every
other layer builds on it.}
\label{fig:architecture}
\end{figure}

\subsection{Qubit lines, gates, and markers}

A qubit-line block establishes one qubit's timeline row and carries no adjustable state of
its own. A single gate block, placed east of a qubit line's anchor, cycles through six
presets on right-click -- Pauli~$X$, $Y$, $Z$, Hadamard, $S$, and $T$ -- exactly as
CircuitSimCraft's diode cycles through its own preset family~\cite{picos2026circuitsimcraft}: one
block class serves every single-qubit gate rather than six separate block classes, keeping
the block/texture/recipe count tractable while still exposing the full single-qubit gate set
this mod supports. A control marker and a target marker, placed at the same east-offset on
two (for CNOT) or three (for Toffoli, with two control markers) different qubit rows,
together represent a multi-qubit gate; unlike every two-terminal component in
CircuitSimCraft, none of these blocks is directional, since a gate's effect does not depend on
which horizontal face a player was looking toward when it was placed.

\subsection{The world-scan-to-circuit bridge}

Unlike CircuitSimCraft's continuous, per-tick network rebuild, QuantumCraft constructs a circuit
only when a player right-clicks the probe against a qubit line, via a four-step scan.
First, a flood fill over six-face adjacency, following only positions whose block entity is a
qubit line, finds every qubit line connected to the clicked position. Second, the found lines
are sorted by position (their north--south coordinate primarily) to assign each a stable qubit
index from~0 up to the three-qubit cap this mod supports. Third, each qubit row is scanned
block by block to the east, up to a fixed maximum timeline length, recording whatever gate or
marker block entity, if any, occupies each offset. Fourth, the recorded entries are grouped by
their shared east-offset (their time column) and translated into circuit operations: a lone
gate entry becomes a single-qubit gate applied to its own row; one control marker paired with
one target marker at the same column becomes a controlled-NOT; two control markers paired
with one target becomes a Toffoli gate. Any other combination found within a single column --
an unpaired control with no target, two target markers, or more controls than the mod's
gate set supports -- is recorded as a warning and simply skipped rather than causing the scan
to fail, so an incompletely wired column fails safely, in the same spirit as CircuitSimCraft's
floating-node handling~\cite{picos2026circuitsimcraft}, rather than crashing or silently
misinterpreting the player's intent.

\subsection{The quantum probe}

A single handheld probe item right-clicked against a qubit line triggers the scan described
above, builds the resulting circuit's final statevector (Section~\ref{sec:quantum-simulator}),
and reports one of two things back to the clicking player, depending on whether the
interaction was a plain right-click or a shift-right-click. A plain right-click performs an
\emph{ideal readout}: every basis state's exact probability is computed and sent to the
client, with no randomness and no collapse, the same sense in which CircuitSimCraft's
oscilloscope displays an exact voltage rather than a noisy measurement. A shift-right-click
instead performs a genuine \emph{measurement}: the final statevector is sampled once, using
the standard Born-rule cumulative-probability procedure (Section~\ref{sec:quantum-simulator}),
the state is collapsed onto that single sampled outcome, and only that one classical outcome
-- not the full probability distribution that produced it -- is reported. This distinction is
made visually unmistakable in the result heads-up display (Figure~\ref{fig:hud}), which
labels the two modes ``IDEAL READOUT'' and ``MEASURED'' respectively and highlights only the
single collapsed row in the latter case, so that the difference between reading out amplitude
information no physical quantum computer could ever expose directly and performing an actual
one-shot measurement is never conflated.
